\subsubsection{The binary stopping game and its lower envelope}
\label{sec:optional-lower}

The first four branches already occur on binary instances
$p_j\in\{0,u\}$.  The stopping point must be hidden from the algorithm;
announcing it in advance gives a strictly weaker lower bound.

\begin{lemma}[Raw-safe hidden stopping]
\label{lem:opt:hidden-stopping}
Fix $u>1$, a target $R\le u$, and a stopping fraction $\alpha\in(0,1)$.
Against a deterministic algorithm, there is a legal binary adversary whose
leading excess at its first crossing is
\begin{equation}
 (u-R)(1-\sigma^2)+\sigma^2 f_{u,R}(y,b),
 \label{eq:opt:stopping-decomposition}
\end{equation}
where $0\le b\le y\le1$ and
\begin{equation}
 0\le y-\alpha<\frac1{n-v}.
 \label{eq:opt:stopping-overshoot}
\end{equation}
Moreover,
\begin{equation}
 f_{u,R}(y,b)=
 1-R+2y+(u-1)(1-R)y^2
 +b^2+2b(u-1-uy).
 \label{eq:opt:stopping-f}
\end{equation}
Here $\sigma\in[0,1]$ is the fraction of jobs not already executed raw.
Consequently, if
\[
 \min_{0\le b\le\alpha}f_{u,R}(\alpha,b)\ge0,
\]
then the adversary gives ratio at least $R-O_u(1/n)$.
\end{lemma}

\begin{proof}
Start in a long phase.  A fresh job that is tested receives $p=u$, while a
fresh job run raw is assigned the hidden value zero.  If $v$ jobs have been
run raw and $L$ tested-long jobs have been seen, stop at the first touch for
which
\begin{equation}
 L\ge\alpha(n-v).
 \label{eq:opt:stopping-line}
\end{equation}
All untouched jobs are then assigned value zero.  Only for the purpose of
lower-bounding the remaining cost, we grant the algorithm this future
information and allow an optimal continuation; in the actual interaction,
values are still revealed only by tests.  If the line is never crossed, every job
was run raw: if even one job had been tested, then after the last first touch
the tested fraction would be one.  On the all-zero hidden instance the ratio
is $u\ge R$.

At a crossing, let $e$ of the $L$ long jobs already be complete, let $d=L-e$
be pending, and let $x=n-v-e-d$ jobs remain untouched.  The following
favorable canonical schedule lower-bounds the algorithm: run the $v$ raw
jobs; test and immediately process the $e$ completed long jobs; perform the
$d$ noncompleting long tests; test the $x$ zero jobs; and finally process the
$d$ pending long jobs.  Indeed, if a noncompleting unit test is immediately
before an available length-$u$ completion while $h$ jobs remain, the two
orders contribute $h+hu$ and $hu+(h-1)$, respectively; moving the
completion left saves one.  Raw completions are available from time zero,
and identical long jobs may be relabeled so that the $e$ completed ones are
tested first.  After the switch, a zero test-completion has length one and
therefore precedes a pending operation of length $u$.  These exchanges
respect precedence.

The exact costs of this relaxed schedule and of OPT are
\begin{align}
 A_{\rm ex}={}&
 u\left(vn-\frac{v(v-1)}2\right)
 +(1+u)\left(e(n-v)-\frac{e(e-1)}2\right) \notag\\
 &+d(n-v-e)+x(d+x)-\frac{x(x-1)}2
 +u\frac{d(d+1)}2,
 \label{eq:opt:stopping-Aexact}\\
 O_{\rm ex}={}&\frac{n(n+1)}2
 +(u-1)\frac{(e+d)(e+d+1)}2.
 \label{eq:opt:stopping-Oexact}
\end{align}
Put $\nu=v/n$, $\delta=(v+e+d)/n$, and $\lambda=e/n$.  Twice the quadratic
coefficients in \cref{eq:opt:stopping-Aexact,eq:opt:stopping-Oexact} are
\begin{align*}
 \mathcal A={}&1+2\delta(1-u\nu)+(u-1)\delta^2
 +2\nu(\nu+u-2)\\
 &+\lambda^2+2\lambda(\nu+u-1-u\delta),\\
 \mathcal O={}&1+(u-1)(\delta-\nu)^2.
\end{align*}
Now put $\sigma=1-\nu$, $y=(e+d)/(n-v)$, and $b=e/(n-v)$.
Direct expansion gives exactly \cref{eq:opt:stopping-decomposition} after
normalizing by $n^2/2$.  The first-crossing overshoot in
\cref{eq:opt:stopping-line} is less than one: a test increases
$e+d-\alpha(n-v)$ by one and a raw first touch increases it by $\alpha$.
This proves \cref{eq:opt:stopping-overshoot}, and in particular
$\sigma(y-\alpha)<1/n$.  There is one small domain issue: the actual
$b\le y$ can exceed $\alpha$.  Put $\bar b=\min\{b,\alpha\}$.  Then
$0\le b-\bar b\le y-\alpha$.  Since $f_{u,R}$ is Lipschitz on the unit
square,
\begin{equation}
 \sigma^2 f_{u,R}(y,b)
 \ge \sigma^2 f_{u,R}(\alpha,\bar b)-O_u(1/n).
\end{equation}
Thus the stated minimum over $0\le\bar b\le\alpha$ suffices uniformly even
when $n-v$ is small.  The discarded diagonal terms are $O_u(n)$ in the
unnormalized costs.
Finally, the adaptive transcript defines a fixed instance by replay: every
tested answer was fixed when revealed, and every raw job's zero value stayed
hidden.  This proves the lemma.
\end{proof}

\begin{proposition}[Binary lower curve]
\label{prop:opt:binary-lower}
Every deterministic algorithm has size-asymptotic ratio at least
\begin{equation}
 B(u)=
 \begin{cases}
  1,&0<u\le1,\\
  u,&\uDet{1}\le u\le \uDet{2},\\
  \Rquad(u),&\uDet{2}\le u\le \uDet{3},\\[1mm]
  \displaystyle1+\frac1{\sqrt{u-1}},&u\ge \uDet{3}.
 \end{cases}
 \label{eq:opt:binary-curve}
\end{equation}
The lower bound uses only $p_j\in\{0,u\}$.
\end{proposition}

\begin{proof}
\medskip
\noindent\textbf{Claim 1 (Two quadratic branches).}
Minimizing \cref{eq:opt:stopping-f} over the algorithm's completion choice
gives
\begin{equation}
 b_*(y)=\max\{0,uy-(u-1)\}.
 \label{eq:opt:bstar}
\end{equation}
The upper constraint $b\le y$ never binds, because
$uy-(u-1)\le y$ for $y\le1$.
There are therefore two elementary quadratic branches.  On the positive
branch, put
\begin{equation}
 A_u=u^2-u+1,
 \qquad
 D_{u,R}=A_u+(u-1)R.
\end{equation}
Then
\begin{equation}
 \min_b f_{u,R}(y,b)
 =1-R-(u-1)^2+2A_uy-D_{u,R}y^2.
 \label{eq:opt:stopping-positive-branch}
\end{equation}
Its maximizer is $\alpha=A_u/D_{u,R}$, and tangency to zero is
\begin{equation}
 \bigl(1-R-(u-1)^2\bigr)D_{u,R}+A_u^2=0.
 \label{eq:opt:stopping-tangency}
\end{equation}
Expanding \cref{eq:opt:stopping-tangency} gives
\cref{eq:opt:Rquad-equation}.  On the zero-$b$ branch,
\begin{equation}
 \min_b f_{u,R}(y,b)
 =1-R+2y+(u-1)(1-R)y^2,
\end{equation}
and tangency gives
\begin{equation}
 R=1+\frac1{\sqrt{u-1}},
 \qquad
 \alpha=\frac1{\sqrt{u-1}}.
 \label{eq:opt:zero-b-tangency}
\end{equation}
The latter point belongs to the zero-$b$ branch iff
$(u-1)^3\ge u^2$, with equality at $u=\uDet{3}$.

\medskip
\noindent\textbf{Claim 2 (The first transition point).}
For $R=u$, the maximum in
\cref{eq:opt:stopping-positive-branch} is nonnegative exactly while
\begin{equation}
 [u(u-1)]^2-u(u-1)-1\le0,
\end{equation}
whose positive equality is $u=\uDet{2}$.

\medskip
\noindent\textbf{Claim 3 (Attainment of the envelope).}
For $\uDet{1}<u\le \uDet{2}$, use $R=u$ and the maximizing point of
\cref{eq:opt:stopping-positive-branch}.  For
$\uDet{2}\le u\le \uDet{3}$, use the positive root $\Rquad(u)$ and
$\alpha=A_u/D_{u,\Rquad(u)}$; the tangency identity makes the maximum zero.
The minimizing completion mass is positive in the interior and becomes zero
at $u=\uDet{3}$.  For $u\ge \uDet{3}$, use the zero-$b$ tangency point from
\cref{eq:opt:zero-b-tangency}.  The branch conditions are precisely the two
inequalities checked above.  Finally, $\RdetRO(u)\ge1$ is trivial for $u\le1$.
In each nontrivial case \cref{lem:opt:hidden-stopping} supplies fixed binary
instances by replay.
\end{proof}


\subsubsection{The cap-free harmonic core}
\label{sec:lower}

This subsection develops the cap-free construction used by the last two
branches of the revealing curve.  Its notation is local to this subsection.
We first record its extremal specialization and then extract the scaled form
reused by the mixed branch.  The construction
has many equally large groups of positive jobs, revealed from longest to
shortest, followed by one large group of zero jobs.  The only design choice is
the gap between two consecutive positive processing times.  Choosing these
gaps harmonically will make all decisions of the online algorithm disappear
from the leading term.

\begin{theorem}[Extremal harmonic core]
\label{thm:lower}
For every deterministic online algorithm $\mathcal B$ and every $\eps>0$,
there are arbitrarily large fixed instances $I$ with unit tests and rational
processing times such that
\[
  \ALG_{\mathcal B}(I)\ge (\RdetOT-\eps)\OPT(I).
\]
\end{theorem}

\begin{proof}
\medskip
\noindent\textbf{Claim 1 (Fixed-input revelation experiment).}
Fix an integer $K$, a rational number $\xi>0$, and a rational slack
$\gamma\in(0,1)$; one may keep $\gamma=1/2$ throughout.  For a large scaling
integer $H$, the instance will contain $H$ jobs of each of $K$ positive types
$0,1,\ldots,K-1$ and $\xi H$ zero jobs.  We choose $H$ to clear all
denominators.

The adversary assigns processing times according to the order in which tests
are started.  The first $H$ tested jobs receive processing time $p_0$, the
next $H$ receive $p_1$, and so on; after the $K$ positive blocks, every
remaining job receives processing time zero.  Thus
\[
  p_0>p_1>\cdots>p_{K-1}>1,
\]
so the algorithm sees the positive types from longest to shortest and learns
about the zero jobs only at the end.

This adaptive description is only a convenient way to construct the input.
After running the interaction against a fixed deterministic algorithm, assign
to every job label the value revealed when that label was tested.  Replaying
the algorithm on this fixed assignment produces exactly the same transcript.
If the algorithm never tests or completes some job, its cost is infinite and
the lower bound is immediate; otherwise every job label is recorded.
Hence it is enough to lower-bound the cost in the revelation experiment.

\medskip
\noindent\textbf{Claim 2 (Shortest-first phase relaxation).}
Let $T_\ell$ be the start of the first test in positive block $\ell$, and let
$T_K$ be the start of the first test in the zero block.  Phase $\ell$ contains
the completions in $(T_\ell,T_{\ell+1}]$; a completion exactly at $T_\ell$
belongs to the preceding phase.  Each $T_\ell$ is the boundary between two
nonpreemptive operations.  Consequently, a previously tested job that is
unfinished at $T_\ell$ has not started processing and still needs its entire
processing time.

We repeatedly use the following elementary observation.  Suppose jobs with
remaining work $v_1,\ldots,v_m$ complete after some time $T$.  If their
completion order is $q_1,\ldots,q_m$, then the $g$th relative completion time
is at least $\sum_{h\le g}v_{q_h}$.  Summing these prefix inequalities shows
that their total relative completion time is minimized by arranging the
$v_i$ in nondecreasing order.  This is merely a lower bound; the online
algorithm need not actually use that order.

For normalized group masses $q_1,\ldots,q_m$ and remaining lengths
$v_1\le\cdots\le v_m$, this shortest-first bound has leading coefficient
\begin{equation}
 \mathcal S((q_g,v_g)_{g=1}^m)
 =\sum_{g=1}^m
   \left(\frac{v_gq_g^2}{2}
    +q_g\sum_{h<g}v_hq_h\right).
 \label{eq:lower-S}
\end{equation}
More precisely, the corresponding contribution for groups of size $q_gH$
is $H^2\mathcal S+O(H)$.  The quadratic term is simply the familiar sum
$v_g(1+\cdots+q_gH)$ inside one group; the cross term is the waiting caused
by all shorter groups.

We ensure that
\begin{equation}
  p_0-p_{K-1}<1.
  \label{eq:lower-span}
\end{equation}
In phase $\ell$, an unfinished job of an older type $j<\ell$ therefore has
remaining length $p_j$, whereas a type-$\ell$ job completed in its own phase
needs both its test and processing, of total length $1+p_\ell$.  Their order
in the shortest-first relaxation is
\[
  p_{\ell-1}<p_{\ell-2}<\cdots<p_0<1+p_\ell.
\]
At $T_K$, an untested zero job has remaining length $1$, so the relaxed tail
starts with all zero jobs and then uses the positive types from shortest to
longest.  These are the only two places where
\cref{eq:lower-span} is used.

\medskip
\noindent\textbf{Claim 3 (Finite accounting identity).}
We now record the online algorithm's possible behavior.  Normalize all job
counts by $H$.  Let
\begin{itemize}
\item $s_j\in[0,1]$ be the mass of type-$j$ jobs completed in their own
      phase;
\item $d_{j,\ell}\ge0$ be the mass of type-$j$ jobs completed in the later
      phase $\ell>j$; and
\item
\[
  X_{j,\ell}=s_j+\sum_{h=j+1}^{\ell-1}d_{j,h}\le1
  \qquad(\ell>j)
\]
be the mass of type $j$ completed before phase $\ell$.  Thus
$X_{j,\ell+1}$ is the mass completed by the end of phase $\ell$.
\end{itemize}
At time $T_\ell$, the machine has already performed the tests of the $\ell H$
jobs in earlier positive blocks and has processed every earlier job that has
completed.  Thus
\begin{equation}
  \frac{T_\ell}{H}
  \ge \ell+\sum_{j<\ell}p_jX_{j,\ell}.
  \label{eq:lower-start}
\end{equation}
For each phase, charge its completions their common start time from
\cref{eq:lower-start}, and charge their relative completion times using
\cref{eq:lower-S}.  Do the same at $T_K$ for the remaining positive jobs
and the $\xi H$ zero jobs.  Expanding and collecting equal variables gives
\begin{align}
 \frac{\ALG}{H^2}
 \ge{}& C_K
 +\sum_{j=0}^{K-1}\left(\frac{s_j^2}{2}-\theta_js_j\right)
 \nonumber\\
 &+\sum_{j<\ell}d_{j,\ell}
 \left[
   \ell-j-\theta_j
   -\sum_{j<k<\ell}(p_j-p_k)X_{k,\ell+1}
 \right]-O(1/H),
 \label{eq:lower-accounting}
\end{align}
where $K,\xi,\gamma$ are fixed while $H\to\infty$, and
\begin{align}
 \theta_j
   &=1-\xi(p_j-1)-\sum_{k>j}(p_j-p_k-1),
   \label{eq:lower-theta}\\
 C_K
   &=\frac{\xi^2}{2}+2\xi K+K^2+P_K,
 &
 P_K&=\sum_{j=0}^{K-1}\left(j+\frac12\right)p_j.
 \label{eq:lower-CP}
\end{align}

For completeness, here is how to audit
\cref{eq:lower-accounting} without guessing its origin.  In phase $\ell$,
the masses in the shortest-first relaxation are
\[
 d_{\ell-1,\ell},d_{\ell-2,\ell},\ldots,d_{0,\ell},s_\ell
\]
with respective lengths
\[
 p_{\ell-1},p_{\ell-2},\ldots,p_0,1+p_\ell.
\]
Multiply their total mass by the right-hand side of
\cref{eq:lower-start} and add the functional
\cref{eq:lower-S}.  In the final phase use masses
\[
 \xi,\ 1-X_{K-1,K},\ldots,1-X_{0,K}
\]
and lengths $1,p_{K-1},\ldots,p_0$.  The constant terms sum to $C_K$.
The terms containing only $s_j$ are $s_j^2/2-\theta_js_j$; the coefficient
of each $d_{j,\ell}$ is precisely the bracket in
\cref{eq:lower-accounting}.  Thus the displayed identity is just the
phasewise shortest-first bound with the phase start times retained.  No
optimality assumption about the online schedule is hidden in it.

The meaning of the expression is useful.  Completing an $s_j$ fraction of a
type in its own phase creates a quadratic completion cost but saves later
waiting cost, producing $s_j^2/2-\theta_js_j$.  A completion postponed to a
later positive phase is represented by $d_{j,\ell}$ and its bracket.  We next
choose the processing times so that every such bracket is nonnegative and all
$\theta_j$ are equal.

\medskip
\noindent\textbf{Claim 4 (Harmonic cancellation).}
Set
\begin{equation}
  p_{K-1}=1+\frac{\gamma}{\xi},
  \qquad
  p_i=p_{i+1}+\frac1{\xi+K-1-i}
  \quad(0\le i<K-1).
  \label{eq:lower-harmonic}
\end{equation}
In other words, measured upward from the shortest type, the consecutive gaps
are
\[
  \frac1{\xi+1},\frac1{\xi+2},\ldots,\frac1{\xi+K-1}.
\]
This is the key design choice.  If $m=K-1-j$, then
\begin{align*}
 p_j-1
   &=\frac\gamma \xi+\sum_{h=1}^m\frac1{\xi+h},\\
 \sum_{k>j}(p_j-p_k)
   &=\sum_{h=1}^m\frac{h}{\xi+h}.
\end{align*}
Adding the first line multiplied by $\xi$ to the second line makes each
denominator cancel:
\begin{equation}
 \xi(p_j-1)+\sum_{k>j}(p_j-p_k)
 =\gamma+\sum_{h=1}^m\frac{\xi+h}{\xi+h}
 =\gamma+m.
 \label{eq:lower-telescope}
\end{equation}
Substitution in \cref{eq:lower-theta} gives the same value for every type,
\begin{equation}
  \theta_j=1-\gamma.
  \label{eq:lower-theta-constant}
\end{equation}
This is why the increments in \cref{eq:lower-harmonic} are harmonic: the
factor $\xi$ coming from the zero block and the multiplicity $h$ coming from the
$h$ shorter positive types add up to the denominator $\xi+h$.

The delayed-completion terms are now harmless as well.  Since
$X_{k,\ell+1}\le1$, their brackets are at least
\begin{align*}
 &\ell-j-(1-\gamma)
   -\sum_{j<k<\ell}(p_j-p_k)\\
 &\hspace{20mm}
 =\gamma+\sum_{j<k<\ell}\bigl(1-(p_j-p_k)\bigr)>0,
\end{align*}
where the final inequality follows from
\cref{eq:lower-span}.  We may therefore drop every term containing
$d_{j,\ell}$.  The remaining minimization is not a multivariable optimization
problem: complete one square for each $j$,
\[
  \frac{s_j^2}{2}-(1-\gamma)s_j
  =\frac12\bigl(s_j-(1-\gamma)\bigr)^2
   -\frac12(1-\gamma)^2.
\]
Because $1-\gamma\in(0,1)$, this minimum is feasible for
$s_j\in[0,1]$.  Hence every online algorithm obeys
\begin{equation}
 \liminf_{H\to\infty}\frac{\ALG}{H^2}
 \ge Q_K
 \defeq C_K-\frac K2(1-\gamma)^2.
 \label{eq:lower-QK}
\end{equation}

It remains to compute the offline benchmark.  Knowing all processing times,
the optimum runs the zero jobs first and the positive types from shortest to
longest.  Equivalently, the pairwise contribution of two jobs is
$1+\min\{p_i,p_j\}$.  Its leading coefficient is
\begin{equation}
  D_K=\frac{\xi^2}{2}+\xi K+\frac{K^2}{2}+P_K,
  \qquad
  \lim_{H\to\infty}\frac{\OPT}{H^2}=D_K.
  \label{eq:lower-DK}
\end{equation}
The four terms respectively account for zero--zero pairs, zero--positive
pairs, the unit-test part of positive--positive pairs, and the processing
part of positive--positive pairs.  Thus the finite-$K$ lower-bound
certificate is simply $Q_K/D_K$.

\medskip
\noindent\textbf{Claim 5 (Smooth limit and its maximum).}
Write
\[
  \alpha=\frac K\xi
\]
for the ratio of positive mass to zero mass, and keep $\alpha$ fixed while
$K\to\infty$.  We restrict to $0<\alpha<e-1$.  Then
\[
 p_0-p_{K-1}
 =\sum_{h=1}^{K-1}\frac1{\xi+h}
 <\ln\!\left(1+\frac K\xi\right)
 =\ln(1+\alpha)<1,
\]
which verifies \cref{eq:lower-span}.

The processing contribution $P_K$ has the exact form
\begin{equation}
 P_K
 =\frac{K^2}{2}\left(1+\frac\gamma \xi\right)
  +\frac12\sum_{h=1}^{K-1}\frac{(K-h)^2}{\xi+h}.
 \label{eq:lower-P-exact}
\end{equation}
Indeed, the common baseline $1+\gamma/\xi$ contributes the first term, and
swapping the order of summation shows that the increment $1/(\xi+h)$ receives
total weight $(K-h)^2/2$.  After division by $\xi^2$, the sum becomes a
Riemann integral:
\begin{align}
 \frac{P_K}{\xi^2}&\longrightarrow I(\alpha)
   =\frac{\alpha^2}{2}
    +\frac12\int_0^\alpha\frac{(\alpha-t)^2}{1+t}\,dt
   \nonumber\\
 &=\int_0^\alpha(\alpha-t)\bigl(1+\ln(1+t)\bigr)\,dt
   \nonumber\\
 &=\frac12(1+\alpha)^2\ln(1+\alpha)
   -\frac\alpha2-\frac{\alpha^2}{4}.
 \label{eq:lower-I}
\end{align}
The second integral form follows either by integration by parts or by summing
the limiting processing-time curve $1+\ln(1+t)$ against the remaining
positive mass $\alpha-t$.
The subtraction in \cref{eq:lower-QK} is only $O(K)$, whereas
$C_K,D_K=\Theta(K^2)$.  Since $C_K-D_K=\xi K+K^2/2$, we obtain
\begin{equation}
 R(\alpha)
 \defeq\lim_{K\to\infty}\frac{Q_K}{D_K}
 =1+
 \frac{\alpha+\alpha^2/2}
 {\frac12+\frac\alpha2+\frac{\alpha^2}{4}
  +\frac12(1+\alpha)^2\ln(1+\alpha)}.
 \label{eq:lower-Ralpha}
\end{equation}

This last one-variable maximization is where the constant in the theorem
comes from.  Put $z=(1+\alpha)^2$.  Then
\begin{equation}
  R(z)=1+\frac{2(z-1)}{1+z+z\ln z}.
  \label{eq:lower-Rz}
\end{equation}
The derivative of the fraction has the sign of
\[
  3-z+\ln z.
\]
For $z>1$ this expression is strictly decreasing, is positive at $z=1$, and
tends to $-\infty$.  Hence it has one zero $\uDet{5}$, characterized by
\[
  \uDet{5}-3=\ln \uDet{5}.
\]
The derivative is positive before $\uDet{5}$ and negative afterward, so this is
the global maximum.  At the root,
\[
  1+\uDet{5}+\uDet{5}\ln \uDet{5}=(\uDet{5}-1)^2,
\]
and therefore
\[
  R(\uDet{5})=1+\frac{2}{\uDet{5}-1}=\RdetOT.
\]
Moreover, $\sqrt{\uDet{5}}-1=1.12255\ldots<e-1$, so the maximizer lies in the
range where \cref{eq:lower-span} holds.

\medskip
\noindent\textit{Final parameter choice.}
Choose a rational $\alpha$ sufficiently close to
$\alpha_\star=\sqrt{\uDet{5}}-1$, and choose any rational
$\gamma\in(0,1)$.  For a sufficiently large $K$, set $\xi=K/\alpha$ and use
\cref{eq:lower-harmonic}; then $Q_K/D_K>\RdetOT-\eps/2$.
All parameters are rational.  Let $H$ tend to infinity through multiples of
their common denominators.  \Cref{eq:lower-QK,eq:lower-DK} make the actual ratio exceed
$\RdetOT-\eps$ for all sufficiently large $H$.  Finally, the replay argument
at the beginning of the proof converts the revelation experiment into a fixed
instance for the given deterministic algorithm.
\end{proof}

The mixed branch needs the same construction at an arbitrary total mass and
at a prescribed point of its limiting curve.  The next lemma packages this
rescaled form so that no phase accounting has to be repeated later.

\begin{lemma}[Scaled harmonic core]
\label{lem:opt:harmonic-block}
Fix a total normalized mass $s>0$ and a parameter $t\in[1,e)$.  Against
every deterministic obligatory-testing algorithm, there are scales
$M_\nu\to\infty$ and finite rational multilevel instances with
$sM_\nu+o(M_\nu)$ jobs such that
\begin{equation}
 \frac{\OPT}{M_\nu^2}\longrightarrow
 O_0=\frac{s^2}{4}(1+t^{-2}+2\ln t),
 \qquad
 \liminf_{\nu\to\infty}\frac{\ALG}{M_\nu^2}
 \ge A_0^{\rm lb}
 \defeq O_0+\frac{s^2}{2}(1-t^{-2}).
 \label{eq:opt:harmonic-coefficients}
\end{equation}
Their normalized total positive processing mass converges to
\begin{equation}
 \frac1{M_\nu}\sum_{j:p_j>0}p_j\longrightarrow L=s\ln t,
 \label{eq:opt:harmonic-processing-mass}
\end{equation}
and their largest processing time converges to $1+\ln t<2$.
The adaptive revelation is converted by replay into a fixed instance.
\end{lemma}

\begin{proof}
The case $t=1$ is obtained by continuity from instances with vanishing
positive mass, so suppose $1<t<e$.  In the construction above set
$\alpha=t-1=K/\xi$.  Before rescaling, the total normalized mass is
$K+\xi=\xi t$.  Give every unit of normalized mass the common factor
$s/(\xi t)$.

The phase accounting, harmonic cancellation, and completion of squares in
\cref{eq:lower-accounting,eq:lower-telescope,eq:lower-QK} are homogeneous of
degree two in the job masses.  The Riemann-sum limit
\cref{eq:lower-I} substituted into \cref{eq:lower-DK} therefore gives
\[
 O_0=\frac{s^2}{4}(1+t^{-2}+2\ln t).
\]
Likewise $C_K-D_K=\xi K+K^2/2$, while the $O(K)$ square-completion loss in
\cref{eq:lower-QK} disappears after division by the quadratic scale.  After
the same rescaling this gives the guaranteed lower benchmark
\[
 A_0^{\rm lb}-O_0=\frac{s^2}{2}(1-t^{-2}).
\]
Finally, summing the limiting processing-time curve
$1+\ln(1+x)$ over the positive mass gives $L=s\ln t$, and
\cref{eq:lower-harmonic} gives $p_0\to1+\ln t$.  Rational approximation,
clearing denominators, and the replay argument used in the theorem produce a
diagonal sequence of arbitrarily large fixed rational instances with the
stated limit and liminf bounds.
\end{proof}

\subsubsection{The unrestricted mixed lower bound}
\label{sec:opt:mixed-lower}

For $\uDet{4}<u<\uDet{5}$, neither the binary family nor the pure cap-free
family is tight.  The correct adversary puts a capped block before a scaled
harmonic core.  Two lemmas ensure that raw
execution and early processing of caps cannot escape it.

The cap-free construction needed below is exactly the scaled harmonic core
from \cref{lem:opt:harmonic-block}; no second phase-accounting proof is
needed here.  We now prepend capped jobs to that reusable core.

\begin{lemma}[Cap deferral]
\label{lem:opt:cap-deferral}
Prepend capped mass $m$ with value $p=u$ to a later cap-free block of mass
$s=1-m$ and total processing mass $L$.  Suppose every later value satisfies
$1+p\le u$ and
\begin{equation}
 s(u-1)-L-m\ge0.
 \label{eq:opt:cap-deferral-condition}
\end{equation}
Then processing any capped job before all later jobs are complete cannot
decrease the leading online cost.
\end{lemma}

\begin{proof}
A cap completed after its test block but before a later job of value $p$
has mutual cost $1+u$; leaving it pending until after that job gives $2+p$.
The exchange saves $u-1-p\ge0$.  If cap mass $q$ is instead completed inside
the cap-test block, its pairs with all later jobs add
$q[s(u-1)-L]$.  Its largest possible saving against not-yet-tested caps is
$qm-q^2/2$.  Relative to leaving every cap pending, the change is therefore
at least
\begin{equation}
 q[s(u-1)-L-m]+\frac{q^2}{2}\ge0.
\end{equation}
The finite diagonal discrepancy is lower order.
\end{proof}

\begin{lemma}[Capped-block accounting]
\label{lem:opt:mixed-accounting}
After cap deferral, the offline coefficient $O$ and the guaranteed online
lower benchmark $A^{\rm lb}$ satisfy
\begin{align}
 A^{\rm lb}={}&A_0^{\rm lb}+m(1-m)+m(1+L)+\frac u2m^2,
 \label{eq:opt:mixed-A}\\
 O={}&O_0+m(1-m+L)+\frac u2m^2.
 \label{eq:opt:mixed-O}
\end{align}
In particular, $A^{\rm lb}-O=(A_0^{\rm lb}-O_0)+m$.
\end{lemma}

\begin{proof}
A cap--later pair contributes $2+p$ online and $1+p$ offline,
whereas a cap--cap pair contributes $2+u$ online and $u$ offline; collecting
these pairs yields the two cross terms and the cap squares above.
Subtracting the two identities gives the final assertion.
\end{proof}

\begin{lemma}[Raw-safe quota and replay]
\label{lem:opt:mixed-raw-safe}
Fix $m\in(0,1)$ and a later cap-free block satisfying the hypotheses of
\cref{lem:opt:cap-deferral}.  Against every deterministic revealing algorithm
there are fixed instances whose asymptotic excess is at least the cap-free
hybrid benchmark $A^{\rm lb}-R\,O$ from
\cref{eq:opt:mixed-A,eq:opt:mixed-O} for every target
$R\le u$.
\end{lemma}

\begin{proof}
In the initial phase, a tested fresh job receives $p=u$, while a job run raw
is assigned hidden value zero.  If $v$ raw jobs and $\ell$ tested caps have
been touched, stop at the first crossing
\begin{equation}
 \frac{\ell}{n-v}\ge m.
 \label{eq:opt:mixed-quota}
\end{equation}
With $S=n-v$, the overshoot is less than one, so $\ell=mS+O(1)$.  If no
crossing occurs, no test can have occurred: after all first touches, any
positive $\ell$ would make the ratio in \cref{eq:opt:mixed-quota} equal one.
Thus every job was run raw and the all-zero hidden instance has ratio
$u\ge R$.

At a crossing, scale the prescribed later instance to the remaining
$S-\ell$ first touches, absorbing the one-job cap overshoot and integer
rounding into its zero group.  A later test reveals the next value in this
fixed type stream; a later raw action consumes the next value but keeps it
hidden from the revealing algorithm.

Move every initial raw-zero completion to the start.  It has the same length
$u$ as cap processing, and crossing a unit test only lowers completion cost.
With
\begin{equation}
 Z_v=vn-\frac{v(v-1)}2,
\end{equation}
this prefix contributes $uZ_v$ online and $Z_v$ offline.  Delete it.  The
remaining $S$ labels consist, up to one-job rounding, of capped fraction $m$
and the prescribed later cap-free block.

If the revealing algorithm later runs a fresh assigned value $p$ raw, a
physical cap-free simulation tests and immediately processes that job in
time $1+p\le u$, suppresses the answer from the virtual revealing algorithm,
and idles until its virtual $u$-block ends.  Inductively every physical
completion is no later than its virtual revealing completion.  Since the
finite-cap and cap-free offline lengths are both $1+p$ on the later block,
\begin{equation}
 \ALG_{\mathsf{RO}}-R\OPT_{\mathsf{RO}}
 \ge (u-R)Z_v+
 \bigl(\ALG_{\rm hyb}-R\OPT_{\rm hyb}\bigr)-O_u(n).
 \label{eq:opt:mixed-raw-decomposition}
\end{equation}
The first term is nonnegative.  Finally, record every assigned value on its
job label.  Replaying this fixed vector produces the same transcript because
the values learned only by the physical simulation were hidden from the
virtual algorithm.  This proves legality.
\end{proof}

\begin{proposition}[Mixed lower branch]
\label{prop:opt:mixed-lower}
For every $\uDet{4}\le u\le\uDet{5}$, every deterministic revealing
algorithm has size-asymptotic competitive ratio at least $\Rmix(u)$.
Equivalently,
\begin{equation}
 \RdetRO(u)\ge\Rmix(u).
 \label{eq:opt:mixed-lower-final}
\end{equation}
\end{proposition}

\begin{proof}
\medskip
\noindent\textbf{Claim 1 (Two-parameter lower ratio).}
Apply \cref{lem:opt:harmonic-block} with $s=1-m$.  Substitution into
\cref{eq:opt:mixed-A,eq:opt:mixed-O} gives the guaranteed two-parameter
lower ratio
\begin{equation}
 \mathcal R_u(m,t)=1+
 \frac{m+\frac{(1-m)^2}{2}(1-t^{-2})}
 {\frac{(1-m)^2}{4}(1+t^{-2}+2\ln t)
  +m(1-m)(1+\ln t)+\frac u2m^2}.
 \label{eq:opt:mixed-ratio}
\end{equation}
\medskip
\noindent\textbf{Claim 2 (Mixed parametrization).}
At an interior stationary point with value $1+c$, differentiation and the
zero-gap equation give
\begin{equation}
 \ln t=u-2-\frac1c,
 \qquad
 t^2=\frac{1-m}{1+m}\left(1+\frac2c\right).
 \label{eq:opt:mixed-stationarity}
\end{equation}
Eliminating $t$ gives precisely
\cref{eq:opt:mixed-parametrization}.  Conversely, those equations make
$\mathcal R_u(m,t)=1+c$.

\medskip
\noindent\textbf{Claim 3 (Feasibility and uniqueness).}
The branch runs from
\begin{equation}
 (u,m,t)=\left(\uDet{4},\frac1\varphi,1\right)
 \quad\text{to}\quad
 (u,m,t)=(\uDet{5},0,\sqrt{\uDet{5}}).
 \label{eq:opt:mixed-endpoints}
\end{equation}
It satisfies $u\ge\uDet{4}$, $\ln t<1$, and $m\le1/\varphi$.  Therefore
\begin{equation}
 (1-m)(u-1-\ln t)
 >(1-m)\varphi\ge\frac1\varphi\ge m,
\end{equation}
so \cref{eq:opt:cap-deferral-condition} holds; also $1+p<3<u$ for every
later positive value in a sufficiently fine finite construction.  The
functions
\begin{equation}
 g(c)=1-\frac1c+\frac12\ln\!\left(1+\frac2c\right),
 \qquad h(m)=\operatorname{artanh}m-m
\end{equation}
have $g'(c)=2/[c^2(c+2)]>0$ and
$h'(m)=m^2/(1-m^2)>0$.  Hence $m(c)=h^{-1}(g(c))$ is strictly increasing,
while
\begin{equation}
 u(c)=1+\frac2c-m(c)
\end{equation}
is strictly decreasing.  The identities $g(r)=h(0)=0$ and
$g(1/\varphi)=h(1/\varphi)$ give the two endpoints in
\cref{eq:opt:mixed-endpoints}; consequently every
$u\in[\uDet{4},\uDet{5}]$ has exactly one parameter pair.

\medskip
\noindent\textbf{Claim 4 (Fixed finite instances).}
To pass to fixed finite instances, first approximate $(m,t)$ rationally
while retaining the strict cap inequalities, then choose a sufficiently
fine harmonic family, and finally multiply all rational block sizes by a
common integer $H$.  For fixed harmonic depth, quota overshoot, diagonals,
and rounding are $O_u(H)$ whereas the displayed costs are $\Theta(H^2)$.
Send first $H$, then the harmonic depth, and finally the rational
approximation to their limits.  Together with
\cref{lem:opt:mixed-raw-safe}, this proves the proposition.
The construction above applies directly in the open interval.  At
$u=\uDet{4}$ and $u=\uDet{5}$ the same statement follows by approaching from
the interior, so the endpoints are included.
\end{proof}

\subsubsection{The plateau lower bound by cap-free transfer}

The harmonic construction contains no capped outcomes.  The next lemma
shows that any revealing algorithm at cap $u$ induces an obligatory-testing
algorithm on inputs with $p_j\le u-1$, so the cap-free lower bound transfers
without changing either the transcript comparison or the offline optimum.

\begin{lemma}[Bounded cap-free transfer]
\label{lem:opt:obligatory-transfer}
Let $u\ge1$, and let $L_{\mathsf{OT}}(P)$ denote the deterministic
size-asymptotic value in the cap-free limit, restricted to
$0\le p_j\le P$.  Then
\begin{equation}
 \RdetRO(u)\ge L_{\mathsf{OT}}(u-1).
 \label{eq:opt:obligatory-transfer}
\end{equation}
In particular,
\begin{equation}
 \RdetRO(u)\ge \RdetOT
 \qquad
 \left(u\ge1+\frac1r=2.752620747896442\ldots\right).
 \label{eq:opt:high-lower}
\end{equation}
\end{lemma}

\begin{proof}
Given a finite-cap algorithm, construct a cap-free algorithm on instances
with $p_j\le u-1$ by running a virtual copy of the finite-cap algorithm.  Copy
ordinary tests and processing actions.  When the virtual algorithm runs an
untouched job raw for $u$ time, physically test and immediately process it in
$1+p_j\le u$ time, hide the answer from the virtual copy, and wait for the
virtual clock.  Every physical completion is no later than its virtual
completion.  Moreover, $\min\{u,1+p_j\}=1+p_j$, so finite-cap and cap-free
optima coincide.  Thus we obtain the following game-value comparison.  For
every finite-cap algorithm
$\mathcal A$, the simulation produces a cap-free algorithm $\mathcal B$ and,
pointwise on every instance with $p_j\le u-1$,
\[
 \ALG_{\mathcal A}\ge\ALG_{\mathcal B},
 \qquad \OPT_u=\OPT_\infty.
\]
Consequently
\[
 \CR_\infty(\mathcal A)
 \ge \CR_\infty(\mathcal B; p\le u-1)
 \ge L_{\mathsf{OT}}(u-1).
\]
Taking the infimum over $\mathcal A$ proves
\cref{eq:opt:obligatory-transfer}.

The harmonic cap-free family attaining $\RdetOT$ has limiting largest
processing time
\begin{equation}
 1+\ln(1+y_*)
 =1+\frac12\ln \uDet{5}
 =\frac{\uDet{5}-1}{2}=\frac1r,
 \qquad y_*=\sqrt{\uDet{5}}-1.
\end{equation}
Approximation from below also gives the endpoint, proving
\cref{eq:opt:high-lower}.
\end{proof}
